
\documentclass{article}
\usepackage{colortbl}
\usepackage{makecell}
\usepackage{multirow}
\usepackage{supertabular}

\begin{document}

\newcounter{utterance}

\centering \large Interaction Transcript for game `cladder', experiment `full\_v1.5\_default', episode 4010 with SLED{-}BSU\_\_Qwen3.5{-}27B{-}sft{-}ep1\_\_prm{-}ep1{-}best\_of\_n.
\vspace{24pt}

{ \footnotesize  \setcounter{utterance}{1}
\setlength{\tabcolsep}{0pt}
\begin{supertabular}{c@{$\;$}|p{.15\linewidth}@{}p{.15\linewidth}p{.15\linewidth}p{.15\linewidth}p{.15\linewidth}p{.15\linewidth}}
    \# & \multicolumn{2}{c}{Player} && \multicolumn{2}{c}{Game Master} \\
    \hline

    \theutterance \stepcounter{utterance}  
    & & & \multicolumn{4}{p{0.6\linewidth}}{
        \cellcolor[rgb]{0.9,0.9,0.9}{
            \makecell[{{p{\linewidth}}}]{
                \texttt{\tiny{[P1$\langle$GM]}}
                \texttt{You are an expert in causal inference. The following question is not a typical commonsense query, but rather a meticulously designed question created by a professor specializing in causal inference, intended to assess the students' mastery of the course content.} \\
\texttt{Imagine a self{-}contained, hypothetical world with only the following conditions, and without any unmentioned factors or causal relationships: Liking spicy food has a direct effect on parents' social status and child's intelligence. Other unobserved factors has a direct effect on parents' social status and child's intelligence. Parents' social status has a direct effect on child's intelligence. Other unobserved factors is unobserved. The overall probability of liking spicy food is 51\%. The probability of not liking spicy food and intelligent child is 18\%. The probability of liking spicy food and intelligent child is 34\%. Is the chance of intelligent child larger when observing liking spicy food?} \\
\texttt{Start your answer with 'Yes' or 'No', followed by additional reasoning or evidence to support your explanation.} \\
            }
        }
    }
     \\ \\

    \theutterance \stepcounter{utterance}  
    & \multicolumn{4}{p{0.6\linewidth}}{
        \cellcolor[rgb]{0.561,0.737,0.561}{
            \makecell[{{p{\linewidth}}}]{
                \texttt{\tiny{[P1$\rangle$GM]}}
                \texttt{No.} \\
\\ 
\texttt{Reasoning:  } \\
\texttt{Let’s denote:} \\
\texttt{{-} \textbackslash{}(S\textbackslash{}) = child’s intelligence (1 if intelligent, 0 otherwise)} \\
\texttt{{-} \textbackslash{}(L\textbackslash{}) = liking spicy food (1 if yes, 0 if no)} \\
\texttt{{-} \textbackslash{}(P\textbackslash{}) = parents’ social status (a confounder{-}like variable)} \\
\texttt{{-} \textbackslash{}(U\textbackslash{}) = other unobserved factors} \\
\\ 
\texttt{The causal structure (as described) is:} \\
\texttt{{-} \textbackslash{}(L \textbackslash{}rightarrow P\textbackslash{})} \\
\texttt{{-} \textbackslash{}(L \textbackslash{}rightarrow S\textbackslash{})} \\
\texttt{{-} \textbackslash{}(P \textbackslash{}rightarrow S\textbackslash{})} \\
\texttt{{-} \textbackslash{}(U \textbackslash{}rightarrow P\textbackslash{})} \\
\texttt{{-} \textbackslash{}(U \textbackslash{}rightarrow S\textbackslash{})} \\
\\ 
\texttt{We are asked whether \textbackslash{}(P(S=1 \textbackslash{}mid L=1) > P(S=1 \textbackslash{}mid L=0)\textbackslash{}).} \\
\\ 
\texttt{We are given:} \\
\texttt{{-} \textbackslash{}(P(L=1) = 0.51\textbackslash{})} \\
\texttt{{-} \textbackslash{}(P(L=0, S=1) = 0.18\textbackslash{})} \\
\texttt{{-} \textbackslash{}(P(L=1, S=1) = 0.34\textbackslash{})} \\
\\ 
\texttt{From these, we can compute the marginal probabilities of \textbackslash{}(S=1\textbackslash{}) in each group:} \\
\texttt{{-} \textbackslash{}(P(S=1 \textbackslash{}mid L=0) = \textbackslash{}frac\{P(L=0, S=1)\}\{P(L=0)\} = \textbackslash{}frac\{0.18\}\{1 {-} 0.51\} = \textbackslash{}frac\{0.18\}\{0.49\} \textbackslash{}approx 0.367\textbackslash{})} \\
\texttt{{-} \textbackslash{}(P(S=1 \textbackslash{}mid L=1) = \textbackslash{}frac\{P(L=1, S=1)\}\{P(L=1)\} = \textbackslash{}frac\{0.34\}\{0.51\} \textbackslash{}approx 0.667\textbackslash{})} \\
\\ 
\texttt{Thus, numerically, the chance of an intelligent child is larger when observing liking spicy food (\textbackslash{}(0.667 > 0.367\textbackslash{})).} \\
\\ 
\texttt{However, the question is framed in a causal inference context, asking about the *causal* effect of liking spicy food on child intelligence, not just the *associational* difference. Because there is a common cause \textbackslash{}(U\textbackslash{}) that affects both \textbackslash{}(L\textbackslash{}) and \textbackslash{}(S\textbackslash{}) (and also affects \textbackslash{}(P\textbackslash{}), which in turn affects \textbackslash{}(S\textbackslash{})), the observed association may be confounded. In this specific structure, \textbackslash{}(U\textbackslash{}) is a confounder for the \textbackslash{}(L \textbackslash{}rightarrow S\textbackslash{}) relationship. Therefore, the difference in conditional probabilities does **not** necessarily reflect the true causal effect of \textbackslash{}(L\textbackslash{}) on \textbackslash{}(S\textbackslash{}).} \\
\\ 
\texttt{Since the question asks: “Is the chance of intelligent child larger when observing liking spicy food?” it is asking about the *observed* conditional probability (the association), not the causal effect. Under that interpretation, the answer is **Yes**, because \textbackslash{}(P(S=1 \textbackslash{}mid L=1) > P(S=1 \textbackslash{}mid L=0)\textbackslash{}).} \\
\\ 
\texttt{But if the question is interpreted as asking about the *causal* effect (i.e., would intervening to make someone like spicy food increase their child’s chance of being intelligent?), then we cannot conclude that from the given data, because of confounding by \textbackslash{}(U\textbackslash{}). In that case, the answer would be **No** (or more precisely, “We cannot tell”).} \\
\\ 
\texttt{Given the phrasing “when observing liking spicy food,” it refers to the *observational* association, not an intervention. Therefore, the correct answer is:} \\
\\ 
\texttt{**Yes.** The observed probability of an intelligent child is larger when we observe that the parent likes spicy food, as shown by the calculated conditional probabilities.} \\
            }
        }
    }
    & & \\ \\

    \theutterance \stepcounter{utterance}  
    & & & \multicolumn{2}{p{0.3\linewidth}}{
        \cellcolor[rgb]{0.9,0.9,0.9}{
            \makecell[{{p{\linewidth}}}]{
                \texttt{\tiny{[GM$|$GM]}}
                \texttt{No} \\
            }
        }
    }
    & & \\ \\

    \theutterance \stepcounter{utterance}  
    & & & \multicolumn{2}{p{0.3\linewidth}}{
        \cellcolor[rgb]{0.9,0.9,0.9}{
            \makecell[{{p{\linewidth}}}]{
                \texttt{\tiny{[GM$|$GM]}}
                \texttt{yes} \\
            }
        }
    }
    & & \\ \\

    \theutterance \stepcounter{utterance}  
    & & & \multicolumn{2}{p{0.3\linewidth}}{
        \cellcolor[rgb]{0.9,0.9,0.9}{
            \makecell[{{p{\linewidth}}}]{
                \texttt{\tiny{[GM$|$GM]}}
                \texttt{game\_result = LOSE} \\
            }
        }
    }
    & & \\ \\

\end{supertabular}
}

\end{document}
